%Recomendaciones
A partir de los hallazgos técnicos, las métricas de rendimiento obtenidas y los paradigmas actuales de la computación en la nube, se plantea directrices estratégicas orientadas a la optimización, escalabilidad y futura puesta en producción del Módulo de Ruteo Seguro. En primer lugar, para erradicar la latencia de red evidenciada durante el cálculo algorítmico, se recomienda imperativamente ejecutar el despliegue del backend (actualmente en Nest.js local) hacia una infraestructura en la nube, esta ubicación  minimizará drásticamente los tiempos de transferencia de paquetes, optimizando el tiempo de respuesta del motor de ruteo y dotando a la plataforma de la alta disponibilidad requerida.

En complemento a la mejora infraestructural, resulta altamente recomendable la integración de una capa de caché espacial en memoria dentro del backend. Dado que el cálculo de rutas seguras y la intersección matemática de polígonos de riesgo demandan un alto costo computacional y de lectura en disco, el almacenamiento temporal de las geometrías correspondientes a los trayectos más solicitados aliviará de manera significativa la sobrecarga de transacciones sobre el motor relacional. Esta estrategia arquitectónica garantizará que el sistema pueda escalar horizontalmente frente a un incremento masivo en la concurrencia de usuarios sin degradar la experiencia de navegación.

Desde una perspectiva algorítmica y funcional, se sugiere evolucionar el modelo matemático de ponderación hacia un enfoque de ruteo dinámico temporal. Considerando que los índices de criminalidad y las dinámicas urbanas presentan fluctuaciones drásticas dependiendo de la franja horaria, la implementación de modelos de aprendizaje automático o algoritmos heurísticos que ajusten los pesos de peligrosidad en tiempo real elevaría exponencialmente la precisión preventiva del aplicativo. Bajo este esquema, el sistema tendría la capacidad adaptativa de ofrecer rutas diferenciadas para un mismo trayecto en horarios diurnos frente a escenarios nocturnos, mitigando los riesgos según la ventana temporal específica.

En lo que concierne con la mejora del frontend se aconseja enriquecer con un sistema de navegación asistida paso a paso. Proveer instrucciones auditivas detalladas y retroalimentación direccional mediante vibraciones hápticas evitará que el ciudadano deba consultar continuamente la pantalla de su dispositivo móvil mientras camina. Esta mejora no solo refinará la accesibilidad y la ergonomía del software, sino que reducirá significativamente las distracciones visuales del peatón, atacando directamente el riesgo inherente de exposición a hurtos de dispositivos electrónicos en la vía pública.